\chapter{Documentação técnica complementar}\label{ap:fontes-tecnicas}

A Tabela~\ref{tab:fontes-tecnicas} apresenta os guias do repositório indicado na seção~\ref{sec:4.6}. Os caminhos partem da raiz do repositório. O guia de arquitetura indica onde consultar a implementação de cada sistema.

\begin{table}[htbp]
\centering
\caption{Guias técnicos do projeto}\label{tab:fontes-tecnicas}
\begingroup
\small
\setlength{\tabcolsep}{4pt}
\renewcommand{\arraystretch}{1.18}
\begin{tabular}{>{\raggedright\arraybackslash}p{0.22\textwidth}>{\raggedright\arraybackslash}p{0.32\textwidth}>{\raggedright\arraybackslash}p{0.37\textwidth}}
\toprule
Guia & Caminho & Conteúdo \\
\midrule
Arquitetura & \nolinkurl{docs/architecture.md} & Cliente, servidor, sincronização e localização do código. \\
Build e execução & \nolinkurl{docs/build.md} & Exportação, partida local e verificação dos artefatos. \\
Configuração & \nolinkurl{docs/configuration.md} & Arquivos, parâmetros e valores padrão. \\
Publicação & \nolinkurl{docs/publishing.md} & Integração contínua, releases e hospedagem. \\
Manutenção & \nolinkurl{docs/maintenance.md} & Mapas, recursos 3D e contratos compartilhados. \\
Métricas & \nolinkurl{docs/metrics.md} & Coleta em CSV e interpretação dos indicadores. \\
\bottomrule
\end{tabular}
\endgroup
\fonte{elaboração própria}
\end{table}

O texto do TCC1 está preservado em \nolinkurl{docs/tcc/TCC1.pdf}.

A síntese do ensaio está em \mbox{\nolinkurl{docs/report/figuras/tcc2/resumo-carga-local.csv}}.
